\section{问题四：波动电价下的严格因果调度}

\subsection{未来价格信息约束与因果性}

若在当天 0:00 直接把附件 4 的完整真实价格轨迹作为未来已知量输入优化，就等价于完全预知未来价格，会低估实际运行成本。为保证策略可实施，本文严格限定每个决策时刻的信息集，并从价格预测层重新构造问题四。

\subsection{因果价格预测}

比较典型均价、1 d 滞后、7 d 滞后、7 d EWMA 和 28 d 滚动中位数后，7 d 滞后同时间隔预测误差最低：
\[
\hat p_{d,t}=p_{d-7,t},\qquad d\ge7,
\tag{14}
\]
启动期采用附件 1 典型日价格。2--12 月预测 MAE 为 \textbf{0.047151 元/kWh}，RMSE 为 0.065348，实际与预测序列相关系数为 \textbf{0.981085}。优化决策使用 $\hat p_{d,t}$，而计划、调整与紧急购电的最终账单全部用附件 4 真实 $p_{d,t}$ 结算。

\subsection{因果模型与完全信息下界}

问题四分别把 Q2 和 Q3 的调度结构替换为预测价格目标函数，得到因果 Q4-2 与因果 Q4-3；同时保留``用真实未来价格优化''的 oracle 作为不可实施的完全信息下界，用于回答``缺少完美价格信息付出了多少成本''。

\begin{longtable}[]{@{}lrrrr@{}}
\toprule
模型 & 2--12 月总费用/元 & 紧急购电量/kWh & Oracle/元 & 因果策略相对下界溢价 \\
\midrule
\endhead
Q4-2 因果 & 15,774,011.240 & 367,486.925 & 15,603,450.520 & 1.093\% \\
Q4-3 因果 & 15,507,241.002 & 261,093.567 & 15,338,096.770 & 1.103\% \\
\bottomrule
\end{longtable}

Q4-3 因果相对 Q4-2 因果节省 \textbf{266,770.238 元（1.691\%）}。在严格消除未来价格泄露后，滚动预报更新的经济收益仍保持，说明这一结论不是由完全预知人为制造。

\begin{figure}[H]
\centering
\includegraphics{../figures_reaslab/fig09_causal_oracle.pdf}
\caption{因果价格策略与完全信息下界}
\end{figure}
